\section{A Structured Note, Model-Free}\label{sec:tour-note}

If any product should break a ledger that refuses to run a pricing model, it is a
structured note: a coupon, a barrier, a payoff computed by proprietary machinery that
takes a quant team to reproduce. This stop runs one through the machine and shows that
the ledger never needs to open that machinery. The specific structure here --- a
periodic coupon with a knock-out barrier --- is an illustrative composition of real,
separately specified mechanisms; the note's price is a black box, and no pricing
formula appears anywhere in this stop, because none ever crosses the door. What crosses
are contractual cash flows and stamped observations, and those are all the ledger ever
needs.

\subsection{The coupon: a scheduled cash move}

The book holds 1{,}000 notes, quoted \EUR{102} cum-coupon, each paying a \EUR{2}
coupon. A coupon date is a dated obligation, so it is a scheduled trigger, and the
lifecycle contract that catches it reads the declared terms and the holding and
submits one transaction. Economically the coupon is the ordinary-dividend boundary
with $(f, c) = (1, 2)$: it books cash $1{,}000 \times 2 = \EUR{2{,}000}$ and advances
the quote from cum to ex, \EUR{102} to \EUR{100}. Value is invariant across the date
up to the cash paid: $1{,}000 \times 102 = \EUR{102{,}000}$ beforehand is
$1{,}000 \times 100 + \EUR{2{,}000}$ after. Because the coupon and the ex-move ride in
one atomic transaction, the two ways it could corrupt --- cash paid while the quote
stays cum, so the coupon is counted twice; or the quote gone ex with no cash, so the
coupon is lost --- are both unreachable. The contract records the coupon as paid in the
unit's status, so a re-fire on the same date pays it once.

\subsection{The barrier: an observation, not an opinion}

The note carries a knock-out barrier. Say, illustratively, the barrier sits at
\EUR{80} on the underlying --- a price-in-basis terms field --- and on the fixing date
the official close prints \EUR{78.40}. The ledger does not decide whether a barrier
``should'' have breached, and it runs no model that would form such a view. The close
enters through the ingest door as a stamped observation, and the contract reads it
against the declared barrier level. If it knocks, the knock is first a recorded change
of the unit's lifecycle state --- ACTIVE to TERMINATED. Whether that termination also
owes cash is a separate matter of the terms: many knock-out structures pay a rebate or
the accrued coupon at termination, and where they do, that leg is a dated settlement
obligation, discharged on its own value date when it falls due, not conjured at the
instant of the breach. The state-change transaction itself is still move-less: the
breach is recorded now, and any cash it triggers settles later as its own obligation.

A move-less transaction may look like a transaction that does nothing, but it is
load-bearing. Without the recorded state change, a later reconstruction of the ledger
as of that date could not tell ``terminated by breach'' from ``still active, and
merely held by no one'' --- two economically distinct states that would otherwise look
identical in the balances. The state change is what the note's changed liability
consists of, so it is recorded even though no cash moved. Consumption of the fixing is
legal only under single-basis consumption: the stamp's basis must agree with the frame
it is read into.

\subsection{Expiry, and the two layers on one number}

At expiry the note settles, and the settlement is where a wrong contract can do its
worst --- and where the ledger's answer to that is exact. Suppose the declared terms
owe \EUR{100} per note, and a defective contract submits a settlement paying \EUR{98}.

The first thing to see is that the door admits it. That \EUR{98} transaction is
balanced, conserving, and basis-consistent; every structural property the door checks
holds. The door has no ground to refuse it, and it should not --- refusing well-formed
transactions on suspicion is exactly how a system of record starts improvising. So the
book now carries $1{,}000 \times 98 = \EUR{98{,}000}$ where \EUR{100{,}000} was owed.
This is the structural layer doing its job: it guarantees the transaction cannot break
the ledger, and it says nothing about whether the economics are right, because it
never depended on the contract being right.

The second thing to see is that the error cannot hide. Everything the settlement was
computed from --- the terms, the adjustment schedule, the registered kinds, the stamped
observations --- is in the log. So the settlement is recomputed from the log prefix and
compared with what was recorded, in canonical form. The recomputation yields \EUR{100},
the record says \EUR{98}, and the comparison reports a difference of \EUR{2} per note,
\EUR{2{,}000} on the 1{,}000-note holding. That nonzero difference is the defect,
located exactly, and it is a defect in the contract, not in the ledger --- the ledger's
part, admitting only well-formed transactions and recording them faithfully, is exactly
the part that was guaranteed. The repair is a later compensating transaction, itself
recorded --- and it is worth separating two things that are easy to run together. That
the error cannot hide is guaranteed by the ledger: the \EUR{2} per note is in the log
for anyone to recompute. That the \EUR{2{,}000} actually comes back is not a ledger
operation at all --- the \EUR{98{,}000} already left the door to an external holder, and
recovering the shortfall is a claims process at the boundary, outside the system of
record. The ledger makes the error undeniable and the correction recordable; it does
not reach across the boundary and unwind a payment already made.

Nothing in this stop asked what the note was worth. The pricing machinery that decided
the coupon, the barrier, and the settlement terms is on the far side of a wall the
ledger never crosses: it records the moves the contract emits and the observations
stamped into the log, and it recomputes those from the log, and that is all. Were an
event to reprice the note, the reprice would enter only as a fresh stamped observation;
the ledger never manufactures a new price by transforming an old one, and any gap
between a transported old price and the next fresh print is left as market profit and
loss, not absorbed into a boundary. \emph{Lifecycle runs on declared terms and stamped
observations alone --- a pricing model's output enters only as a stamped
observation carrying its basis, never as an authoritative number the door interprets or
trusts; the door catches the transactions that are broken, and recomputation from the
log catches the ones that are merely wrong.}
